\section{Kernel: the information lost at zero}
The kernel of a linear map is the set of inputs sent to zero:
\[
\ker T=\{x:T(x)=0\}.
\]
If two inputs have the same output, their difference belongs to the kernel because
\[
T(x-y)=T(x)-T(y)=0.
\]
Thus every collision can be translated into an invisible nonzero input.

\par\noindent\begin{minipage}{\linewidth}
\begin{theorem}\label{thm:kernel}
A linear map is injective if and only if its kernel consists only of the zero vector.
\end{theorem}
\end{minipage}\par

\begin{proof}
If $T$ is injective and $T(x)=0$, then $T(x)=T(0)$. Injectivity gives $x=0$, so the kernel contains no other vector.

Conversely, suppose the only vector mapped to zero is zero. If $T(x)=T(y)$, linearity gives $T(x-y)=0$. The kernel assumption forces $x-y=0$, hence $x=y$. Therefore $T$ is injective.
\end{proof}

The kernel theorem is a two-direction statement about a \emph{linear} map. In its first direction, linearity ensures $T(0)=0$, so an input with output zero can be compared with the zero input. In its second direction, linearity is needed to turn a collision into $T(x-y)=0$. For a nonlinear map, inspecting only inputs sent to zero may miss collisions elsewhere. For instance, $F(x)=x^2$ sends only zero to zero, yet $F(1)=F(-1)$.

This also separates solving a particular equation $T(x)=b$ from proving all targets attainable. A zero kernel says there cannot be two different solutions for any one target. It does not assert that a solution exists for every target. The inclusion $t\mapsto(t,0)$ from $\RR$ to $\RR^2$ is linear and injective, but no input maps to $(0,1)$. To see why counting equations alone cannot establish uniqueness, consider $T(x_1,x_2)=(x_1+x_2,2x_1+2x_2)$. It has two outputs but loses the nonzero vector $(1,-1)$. The second equation contains no new distinction beyond the first.

For the earlier matrix $A$, apply its sum-and-difference rule twice:
\begin{align*}
 A(Ax)&=A(x_1+x_2,x_1-x_2)\\
 &=\big((x_1+x_2)+(x_1-x_2),\ (x_1+x_2)-(x_1-x_2)\big)\\
 &=(2x_1,2x_2)=2x.
\end{align*}
Therefore $Ax=0$ implies $2x=0$, hence $x=0$. The map is injective. Moreover, for any target $b$, taking $x=Ab/2$ gives $Ax=A(Ab)/2=b$. This proves surjectivity and explicitly constructs the inverse. We have proved reversibility without introducing determinants.
